\documentclass[11pt]{article}
\usepackage[T1]{fontenc}
\usepackage{lmodern}
\usepackage{amsmath,amssymb,amsthm,mathtools}
\usepackage{booktabs}
\usepackage{microtype}
\usepackage[hidelinks]{hyperref}
\usepackage[margin=1in]{geometry}
\newtheorem{theorem}{Theorem}
\newtheorem{lemma}{Lemma}
\newtheorem{proposition}{Proposition}
\newcommand{\Li}{\operatorname{Li}}
\newcommand{\rad}{\operatorname{rad}}
\newcommand{\ind}{\mathbf{1}}
\title{An Explicit Sieve Improvement for the Exceptional Exponent of Restrained Odd Perfect Numbers}
\author{Author details to be finalized}
\date{\today}
\begin{document}
\maketitle

\begin{abstract}
A restrained odd perfect number is an odd perfect number having exactly one nonspecial prime factor whose exponent is greater than two. Ochem and Zelinsky proved that this factor must be a power of three and that its exponent \(q\) satisfies \(q\ge 23\,000\,000\,000\).
Building on their structural reduction and small-prime exclusions, we combine a finite congruence sieve with an explicit prime-counting bound for arithmetic progressions and a nonnegative inclusion--exclusion majorant. The resulting computation, expressed using integer and rational arithmetic, gives the improved bound
\[
q>220\,000\,000\,000\,000\,000\,000.
\]
The proof uses a finite set of recorded sieve outputs and a reproducible certificate for the analytic tail estimate. The computational verification status and source dependencies are stated explicitly.
\end{abstract}

\section{Introduction}

For a positive integer \(n\), write \(\sigma(n)\) for the sum of its positive divisors. An odd integer \(N\) is perfect if \(\sigma(N)=2N\). By Euler's theorem, an odd perfect number, if one exists, can be written as
\[
N=p^e\prod_{r}r^{2a_r},\qquad p\equiv e\equiv1\pmod4.
\]
Following \cite{ochem-zelinsky}, call \(N\) \emph{restrained} when precisely one of the even exponents \(2a_r\) exceeds two. Denote that component by \(n^q\).
Ochem and Zelinsky established \(n=3\) and \(q\ge 23\,000\,000\,000\); their proof also excludes primes below \(3\,000\,000\) from the relevant non-\(3\) components and gives a relation between \(q\) and \(\omega(N)\).

Our main result is the following numerical strengthening.

\begin{theorem}\label{thm:main}
If \(N\) is a restrained odd perfect number and \(3^q\) is its unique nonspecial component with exponent greater than two, then
\[
q>220\,000\,000\,000\,000\,000\,000.
\]
\end{theorem}

The calculation has three components. First, a segmented sieve produces an upper bound for the reciprocal sum over possible ordinary prime divisors below \(B=50\,290\,000\,000\,000\). Second, an inclusion--exclusion majorant controls candidates above \(B\). Third, explicit error bounds for primes in arithmetic progressions \cite{bennett-et-al} and a rational verification yield a contradiction for smaller \(q\).

\section{Reduction to a reciprocal-sum inequality}

Use the structural conclusions of \cite{ochem-zelinsky} to write
\[
N=3^q p^e\prod_{r\in R}r^2,\qquad |R|\le 2q+2,
\]
where \(R\) is the set of ordinary prime factors. Every \(r\in R\), as well as the special prime \(p\), exceeds \(3\,000\,000\), according to the cited exclusions.

Since \(\sigma(3^q)/3^q<3/2\), and for a prime \(u\) and \(a\ge1\),
\[
\log\frac{\sigma(u^a)}{u^a}
<\log\frac{u}{u-1}<\frac1{u-1},
\]
we obtain the necessary inequality
\begin{equation}\label{eq:necessary}
\log\frac43<
\frac1{2\,999\,999}+\sum_{r\in R}\frac1{r-1}.
\end{equation}

Fix \(B<Z\), and put
\[
w_Z(r)=\frac1{r-1}-\frac1{Z-1}\qquad (r\le Z).
\]
Using at most \(2q+2\) ordinary prime factors, we will bound the sum on the right of \eqref{eq:necessary} by the finite-sieve contribution, a weighted tail, and a baseline term \((2q+2-K_B)/(Z-1)\).

\section{Finite congruence sieve}\label{sec:finite}

For every ordinary prime divisor \(r\), each prime factor of
\[
\sigma(r^2)=r^2+r+1
\]
must divide \(N\). If \(\ell\equiv1\pmod3\) is excluded as a factor of \(N\), then the two roots of \(X^2+X+1\) modulo \(\ell\) may be removed as possible residue classes for \(r\).

A segmented Eratosthenes sieve removes composites and these forbidden residue classes. For each surviving prime \(r\), the program accumulates the integer upper bound
\[
\left\lceil \frac{10^{18}}{r-1}\right\rceil .
\]
The recorded data are summarized in Table~\ref{tab:sieve}; the complete \(4\,943\)-row table is distributed with the source code.

\begin{table}[ht]
\centering\small
\caption{Finite-sieve output on \((3\,000\,000,B]\), with reciprocal sums scaled by \(10^{18}\).}
\label{tab:sieve}
\begin{tabular}{@{}lrr@{}}
\toprule
Interval & Surviving candidates & Scaled reciprocal upper sum\\
\midrule
\((3\cdot10^6,2\cdot10^{11}]\) & 917\,142\,659 & 62\,696\,136\,715\,648\,189\\
\((2\cdot10^{11},6\cdot10^{11}]\) & 2\,225\,037\,956 & 6\,133\,899\,512\,944\,946\\
\((6\cdot10^{11},B]\) & 186\,184\,186\,178 & 17\,324\,085\,008\,638\,533\\
\midrule
Total & 189\,326\,366\,793 & 86\,154\,121\,237\,231\,668\\
\bottomrule
\end{tabular}
\end{table}

In particular, the recorded values give
\[
K_B=189\,326\,366\,793,\qquad
S_B=0.086154121237231668.
\]
The full finite computation has not yet been independently rerun from start to finish; the published certificate presently checks its saved rows and totals.

\section{A nonnegative sieve majorant}\label{sec:majorant}

Let
\[
P=\{7,13,19,31,37,43,61,67,73,79,97,103\}.
\]
For primes \(r>B\), define
\[
I_\ell(r)=\ind_{\{\ell\mid r^2+r+1\}},\quad
k(r)=\sum_{\ell\in P}I_\ell(r),
\]
and let
\[
T=\{(a,b): a,b\in P\setminus\{7\},\ a<b,\ 7ab\le10^4\}.
\]
Consider the polynomial
\begin{equation}\label{eq:majorant}
G(r)=1-k(r)+\binom{k(r)}2
-I_7(r)\sum_{(a,b)\in T}I_a(r)I_b(r).
\end{equation}

\begin{lemma}\label{lem:majorant}
For every prime \(r>B\),
\[
G(r)\ge \ind_{\{I_\ell(r)=0\ \text{for all }\ell\in P\}}.
\]
\end{lemma}
\begin{proof}
If \(I_7(r)=0\), then \(G=1-k+\binom{k}{2}\), which is \(1\) when \(k=0\) and equals \((k-1)(k-2)/2\ge0\) for \(k\ge1\). If \(I_7(r)=1\), put \(t=k-1\). The value of \(G\) is
\[
\binom{t}{2}-\sum_{(a,b)\in T}I_a(r)I_b(r)\ge0
\]
because the final sum is a subcollection of the \(\binom{t}{2}\) pairs among the \(t\) active indicators other than \(I_7\). The indicator on the right vanishes in the latter case.
\end{proof}

The expansion of \eqref{eq:majorant} has \(97\) monomials, and every nonconstant monomial has squarefree modulus at most \(10^4\). Every modulus is a product of distinct primes \(\ell\equiv1\pmod3\); each nonempty intersection gives \(2^{|S|}\) reduced residue classes modulo \(\prod_{\ell\in S}\ell\). Thus the main coefficient is
\begin{equation}\label{eq:density}
\delta=\sum_{S}a_S\prod_{\ell\in S}\frac{2}{\ell-1}
=\frac{64959667}{169646400}.
\end{equation}

For \(t\ge B\), apply the explicit progression-counting bounds from \cite{bennett-et-al} term by term and take the absolute sum of the errors. Using a modulus-specific coefficient for \(m=1\) and \(m=7\), and the uniform \(1/840\) coefficient for the other moduli, yields
\begin{equation}\label{eq:count}
\left|\sum_{\substack{r\le t\\r\ \mathrm{prime}}}G(r)
-\delta\Li(t)\right|
\le C\frac{t}{(\log t)^2},
\qquad C=\frac{21575999}{42000000}.
\end{equation}
Here the \(m=1\) coefficient is bounded by \(0.0008375\), the per-class \(m=7\) coefficient by \(0.000486\), and the other individual residue-class coefficients by \(1/840\). The hypotheses \(t\ge B\) and \(m\le10^4\) lie inside the respective published ranges. The constants and all \(2^{12}\) Boolean cases are verified by the provided script.

\section{Explicit upper bound for the tail}

Define \(P_G(t)=\sum_{r\le t,\ r\ \mathrm{prime}}G(r)\). Lemma~\ref{lem:majorant} and partial summation give
\begin{equation}\label{eq:tail}
\sum_{\substack{B<r\le Z\\r\in R}}w_Z(r)
\le \int_B^Z\frac{P_G(t)-P_G(B)}{(t-1)^2}\,dt.
\end{equation}
Put \(b=\log B\), \(z=\log Z\), and \(\rho=B/(B-1)\). Inserting \eqref{eq:count} into \eqref{eq:tail}, and bounding the main integral and errors, gives
\begin{align}
M&=\delta\left[\rho\log(z/b)
-\frac{Z/z-B/b+(Z-B)/z^2}{Z-1}\right],\label{eq:M}\\
E&=C\left[\frac{\rho}{b^2}
+\rho^2\left(\frac1b-\frac1z\right)\right].\label{eq:E}
\end{align}
These are upper bounds for the main contribution and the error contribution respectively.

\section{Proof of Theorem~\ref{thm:main}}

Take
\[
B=50\,290\,000\,000\,000,\quad
Z=57\,623\,000\,000\,000\,000\,000\,000,
\]
and suppose \(q\le q_0=220\,000\,000\,000\,000\,000\,000\). Combining the previous estimates gives
\begin{equation}\label{eq:certificate}
\frac1{2\,999\,999}+S_B+M+E+
\frac{2q_0+2-K_B}{Z-1}\le 0.287679827114503
<0.287682072451780\le\log\frac43,
\end{equation}
contradicting \eqref{eq:necessary}. The decimal inequalities in \eqref{eq:certificate} are outward-rounded consequences of rational calculations, using explicit finite series with remainder bounds for logarithms. This proves Theorem~\ref{thm:main}, subject to the finite-sieve exclusions recorded in Section~\ref{sec:finite}. \qed

\section{Computational reproducibility}

The computation and proof certificate are provided in the \texttt{Stage 1} directory of \url{https://github.com/stormwhip/OddPerfectNumberBounds}. The main verification entry point is
\begin{verbatim}
python QBoundVerifyAPRefined.py
\end{verbatim}
which imports the earlier certificate programs, checks the \(4\,943\) saved rows, and performs the final rational comparison. The segmented sieve is implemented in \texttt{QBoundPrefix.cpp}. An independent full rerun of the finite sieve and a separate audit of the cited prime-exclusion tree are pending.

\section*{Acknowledgments}
The structural reductions, factor-tree exclusions, and initial exponent bound used here are due to Pascal Ochem and Joshua Zelinsky. The explicit prime-counting theorem is due to Bennett, Martin, O'Bryant, and Rechnitzer.

\begin{thebibliography}{9}
\bibitem{ochem-zelinsky}
P. Ochem and J. Zelinsky,
\emph{On odd perfect numbers with exactly one even exponent greater than 2},
arXiv:2607.19746 (2026).

\bibitem{bennett-et-al}
M. A. Bennett, G. Martin, K. O'Bryant, and A. Rechnitzer,
\emph{Explicit bounds for primes in arithmetic progressions},
Illinois Journal of Mathematics \textbf{62} (2018), 427--532.
\end{thebibliography}
\end{document}
